\documentclass[12pt]{article}
\usepackage{mathptmx}
\usepackage[T1]{fontenc}
\usepackage{graphicx}
\usepackage[spanish,es-tabla]{babel}
\usepackage{geometry}
\usepackage[table,xcdraw]{xcolor}
\usepackage{colortbl}
\usepackage{multirow}
\usepackage{xurl}
\usepackage{hyperref}
\usepackage{amsmath}
% subcaption conserva el soporte de subfiguras sin conflicto con subfig.
\usepackage{pdflscape}
\usepackage{lscape}
\usepackage{pdfpages}
\usepackage{float}
\usepackage{caption}
\usepackage{subcaption}
\usepackage{fancyhdr}
\usepackage{setspace}
\usepackage{tocloft}
\usepackage{ragged2e}


\renewcommand{\tablename}{Tabla}
\renewcommand{\labelitemi}{\(\bullet\)}

\setlength{\parskip}{1em}%Distancia entre párrafos.
\renewcommand{\baselinestretch}{1.0}% Espacio entre línea y línea o interlineado.





\geometry{margin= 1in}

\setlength{\parindent}{24pt}
\clubpenalty=10000
\widowpenalty=10000

\hypersetup{hidelinks,hypertexnames=false,pdftitle={Anteproyecto de tesis - Matías Arancibia},pdfauthor={Matías Ignacio Arancibia Villanueva}}
\newif\ifpropuestasfiguras
\propuestasfigurastrue % Cambiar a false para ocultar sugerencias editoriales.
\newcommand{\propuestafigura}[2]{\ifpropuestasfiguras
\par\noindent\fcolorbox{gray}{gray!6}{\parbox{0.94\linewidth}{\small\textbf{Propuesta de figura: #1.} #2}}\par
\fi}
\begin{document}

\fancypagestyle{estandar}
{
    \fancyhf{}
    \setlength{\headheight}{30pt}
    \fancyhead[R]{\fontsize{9}{9}\selectfont Universidad Austral de Chile\\28 de septiembre de 2026 \\[-1mm]}
    \fancyhead[L]{\fontsize{13}{13}\selectfont Página \thepage \\[-1mm]}
    \lfoot{} \cfoot{} \rfoot{}
}












\thispagestyle{empty}  % Página sin encabezado ni pie de página

% Logos en la parte superior
\begin{figure}[H]
    \centering
    \includegraphics[width=0.6\linewidth]{Imágenes/uachlogo.png}
\end{figure}
\vspace{0cm}


\begin{flushright}
    \textbf{Profesor Patrocinante:} \\
    Dr. Victor Poblete Ramírez\\
    \textit{Instituto de Acústica} \\
    \textit{Universidad Austral de Chile}




    \textbf{Profesor informante:} \\
    Rodrigo Torres Alarcón\\
    \textit{Instituto de Acústica} \\
    \textit{Universidad Austral de Chile} 

    
\end{flushright}
\vspace{0.5cm}

\begin{center}
    {\LARGE\ Cancelación de feedback acústico en tiempo real: Comparación entre filtros adaptativos y modelos de deep learning bajo condiciones de latencia controlada} \\[0.8 cm]
\end{center}
\vspace{0.3cm}

\begin{flushright}
    \textit{Anteproyecto de tesis para optar al}\\ 
    \textit{Título Profesional de}\\  
    \textit{Ingeniero Civil Acústico }\\
\end{flushright}
\vspace{0.3cm}

 \begin{center}   
    \large Matías Ignacio Arancibia Villanueva \\[0.3cm]
    \large Valdivia, Chile \\[0.3cm]
    \large 28 de septiembre de 2026
\end{center}

\clearpage
\begin{spacing}{0.9}
\tableofcontents
\end{spacing}
\thispagestyle{empty}
\clearpage
\pagestyle{estandar}
\setcounter{page}{1}
\justifying
\section{Introducción}
La realimentación acústica aparece cuando el sonido emitido por un parlante retorna al micrófono y vuelve a amplificarse. En un sistema de audio en vivo, este lazo puede producir coloración, pérdida de claridad y oscilaciones audibles o \textit{howling}. Su control exige conservar la señal deseada y mantener una ganancia útil, sin introducir retardos incompatibles con la aplicación.

Esta investigación comparará dos estrategias de cancelación de realimentación acústica (\textit{Acoustic Feedback Cancellation}, AFC): el método de identificación mediante inyección breve de ruido de Patel y Panahi\cite{patel2020}, denominado ANIA en este trabajo, y el modelo L3C-DeepMFC propuesto por Hao y colaboradores\cite{hao2025}. La abreviatura ANIA se adopta a partir del título \textit{Adaptive Noise Injection Algorithm}; no se atribuye a los autores como sigla oficial. En antecedentes del proyecto también se ha utilizado ASNI para referirse a la inyección adaptativa de ruido de corta duración.

La comparación se desarrollará mediante simulación, verificación de las implementaciones y ensayos físicos en un lazo micrófono--procesador--parlante. Se estudiarán conjuntamente la ganancia adicional estable, la conservación de la voz, la respuesta ante cambios del camino acústico y el costo de ejecución. El banco experimental en C++/JUCE permitirá aplicar un protocolo común y registrar las señales y condiciones de cada ensayo. Posteriormente se evaluará la viabilidad de ejecución en Raspberry Pi 4.

El proyecto ya dispone de infraestructura y verificaciones preliminares. Estas se presentarán como antecedentes de factibilidad; el objetivo del anteproyecto será definir la comparación y las mediciones que se realizarán. No se anticipará la superioridad de ninguno de los métodos.

\clearpage
\section{Planteamiento del Problema}
Una solución de AFC debe reducir el retorno acústico sin confundirlo con la señal que se desea amplificar. En un lazo cerrado, ambas señales están relacionadas, lo que dificulta identificar el camino de realimentación. Además, el desplazamiento del micrófono, la orientación del parlante o la presencia de objetos pueden modificar dicho camino durante la operación.

La identificación con una señal de prueba y la supresión mediante aprendizaje profundo plantean exigencias diferentes. En el primer caso, se deberá considerar la duración y audibilidad de la calibración y la pérdida de eficacia cuando cambia el entorno. En el segundo, será necesario examinar la representatividad de los datos, la ejecución secuencial del modelo y el retardo efectivo de análisis y reconstrucción. Una mejora obtenida sobre archivos procesados fuera de línea no será suficiente para acreditar funcionamiento en un lazo acústico real.

El problema específico de esta tesis será determinar qué compromiso ofrecen ANIA y L3C-DeepMFC entre estabilidad, calidad y recursos bajo condiciones experimentales comparables. Los trabajos seleccionados se desarrollan en contextos de aplicación distintos; por ello, sus resultados publicados no se considerarán directamente transferibles al montaje de esta investigación. Se comprobará su desempeño con señales, caminos acústicos y presupuestos temporales definidos para el banco experimental.

\section{Pregunta de investigación}
¿Qué ventajas y limitaciones presentan ANIA y L3C-DeepMFC para controlar la realimentación acústica de un sistema de audio en vivo, al comparar ganancia adicional estable, calidad de voz, robustez ante cambios del camino acústico y costo computacional bajo restricciones de latencia?

\subsection{Hipótesis de trabajo}
Se plantea que la conveniencia de cada método dependerá del estado del camino acústico y del presupuesto de procesamiento. ANIA podría ofrecer un costo de operación favorable tras identificar un camino aproximadamente estacionario, mientras que L3C-DeepMFC podría conservar mejor la señal en algunas condiciones de realimentación si sus datos de entrenamiento representan dichas condiciones. Esta posible ventaja deberá contrastarse con su latencia, costo y generalización. La hipótesis se evaluará por condición experimental y podrá rechazarse si los datos no muestran estos comportamientos.

\section{Objetivos}
\subsection{Objetivo general}
Comparar experimentalmente ANIA y L3C-DeepMFC para la cancelación de realimentación acústica en tiempo real, considerando ganancia adicional estable, calidad de voz, robustez y costo computacional bajo condiciones de latencia controlada.

\subsection{Objetivos específicos}
\begin{enumerate}
\item Definir un modelo del lazo acústico y un protocolo reproducible que especifique señales, caminos, ganancias, frecuencias de muestreo y condiciones de cambio del entorno.
\item Implementar y verificar la identificación y cancelación de ANIA, distinguiendo la calibración manual del control automático descrito en la publicación original.
\item Implementar, entrenar y verificar L3C-DeepMFC a partir de su descripción publicada, con separación entre entrenamiento, validación y evaluación externa.
\item Integrar ambos métodos en un banco común de procesamiento continuo y comprobar la equivalencia entre las versiones de referencia y las versiones utilizadas en tiempo real.
\item Medir estabilidad, calidad de voz y recuperación ante cambios del camino mediante ensayos emparejados, con incertidumbre y condiciones de validez explícitas.
\item Cuantificar latencia, tiempos de procesamiento y uso de recursos en PC, y evaluar la factibilidad de ejecución sostenida en Raspberry Pi 4.
\end{enumerate}

\section{Justificación}
La contribución esperada será una comparación reproducible que vincule el desempeño acústico con las restricciones de implementación. La ausencia de acople, por sí sola, no acreditará una solución satisfactoria: un método podría evitarlo atenuando excesivamente la señal útil. De manera similar, una mejora de calidad fuera de línea podría perder utilidad si requiere un retardo excesivo o incumple los plazos de procesamiento.

El uso de un banco común permitirá documentar estos compromisos y separar errores del algoritmo de limitaciones del dispositivo o del controlador de audio. La verificación numérica previa reducirá el riesgo de atribuir a un método diferencias causadas por una implementación incorrecta. La evaluación externa con grabaciones propias permitirá examinar condiciones distintas de las utilizadas para ajustar parámetros.

El trabajo aportará implementaciones documentadas, un protocolo de medición y evidencia que permita seleccionar un método según la aplicación. El aporte seguirá siendo válido si el modelo neuronal no supera al método clásico o si no resulta viable en la plataforma embebida, siempre que las conclusiones se sustenten en mediciones comparables.

\section{Marco Teórico}
\subsection{Lazo de realimentación y modelado acústico}
Para el modelo discreto se considerará la señal capturada $d[n]$, la señal deseada en el micrófono $s[n]$, el ruido $v[n]$ y la señal enviada al parlante $u[n]$. Si $h[k]$ representa un camino aproximadamente lineal e invariante durante un intervalo, se tendrá
\begin{equation}
d[n]=s[n]+\sum_{k=0}^{L_h-1}h[k]u[n-k]+v[n].
\label{eq:lazo}
\end{equation}
El modelo incluirá el sonido directo, las reflexiones y los retardos del camino definido entre los puntos de referencia. Se documentará dónde se ubican estos puntos para evitar contabilizar dos veces el retardo de la interfaz o del procesamiento. Cuando cambie el entorno se representará el camino como $h[n,k]$.

En una representación lineal simplificada sin cancelación, si $G(z)$ es el camino de avance, la transferencia entre la señal deseada y la señal capturada contiene el denominador $1-G(z)H(z)$. La estabilidad depende de los polos del lazo cerrado. La condición $|G(e^{j\omega})H(e^{j\omega})|<1$ en todas las frecuencias constituye una condición suficiente bajo las hipótesis usuales de estabilidad de los bloques; el módulo por sí solo no describe todos los casos. En el montaje real se utilizará un criterio operacional de estabilidad, especialmente para métodos no lineales.

\propuestafigura{Lazo acústico y puntos de medición}{Elaborar un diagrama propio con $s[n]$, $d[n]$, $u[n]$, $h[n,k]$, procesador, ganancia y retardos. Marcar la referencia utilizada por el cancelador y los canales exportados. Ubicarlo después de la ecuación~\ref{eq:lazo}.}

\subsection{Identificación y cancelación mediante filtros}
Un cancelador por identificación estimará el retorno mediante $\widehat h[k]$ y lo sustraerá de la entrada:
\begin{equation}
e[n]=d[n]-\sum_{k=0}^{L_{\widehat h}-1}\widehat h[k]u[n-k].
\end{equation}
En el caso ideal, el residuo contendrá la señal deseada y el ruido. Los errores de modelado y los cambios del camino dejarán una componente de realimentación residual. La longitud del filtro, la estimación del retardo y el condicionamiento de la identificación serán variables de diseño relevantes.

Los algoritmos LMS y NLMS constituyen antecedentes de adaptación continua. En esta tesis no se utilizará NLMS como denominación del método de Patel: la versión implementada identifica un filtro mediante una sonda y mantiene sus coeficientes durante la cancelación. Esta diferencia será explícita al describir las variantes experimentales.

\subsection{Método ANIA}
Patel y Panahi proponen identificar el camino acústico mediante una inyección breve de ruido blanco y utilizar el filtro obtenido para cancelar el retorno. Su procedimiento relaciona autocorrelaciones de la sonda y correlaciones cruzadas con la captura. El artículo también contempla detectar el inicio del acople para actualizar el modelo\cite{patel2020}.

En la implementación de esta tesis se verificarán por separado estimación, validación del filtro y operación continua. Se documentarán la regularización numérica, el retardo separado y la validación con una segunda sonda como decisiones de implementación. La modalidad manual servirá para estudiar la cancelación en condiciones controladas; cualquier modalidad automática requerirá validar previamente el detector y se informará como una variante diferenciada.

\subsection{Aprendizaje profundo y L3C-DeepMFC}
L3C-DeepMFC se orienta a estimar voz libre de realimentación mediante procesamiento espectral complejo. Combina modelado recurrente de banda completa y subbandas, reconstrucción por solapamiento y suma de baja latencia, y ajuste fino en lazo cerrado\cite{hao2025}. La arquitectura y el procedimiento de entrenamiento constituirán la referencia para la implementación neuronal de esta tesis.

Se distinguirá entre reproducir el método publicado y desarrollar una variante inspirada en él. Los candidatos MLP, GRU y la red compacta disponibles en el laboratorio del proyecto servirán como herramientas exploratorias; no se presentarán como una reproducción validada de L3C-DeepMFC. Cualquier modificación necesaria para el montaje o la tasa de muestreo se describirá y evaluará explícitamente.

\propuestafigura{Comparación de los dos métodos}{Dibujar dos diagramas propios: calibración y FIR para ANIA; análisis, red recurrente y síntesis para L3C-DeepMFC. Investigar las figuras de las referencias~\cite{patel2020,hao2025} y citar la adaptación. Incluir una línea temporal que muestre las diferencias entre calibración, operación e inferencia.}

\subsection{Tiempo real y latencia}
Para un bloque de $B$ muestras a frecuencia $f_s$, el intervalo disponible es $T_B=B/f_s$. Se distinguirán el tiempo de ejecución de cada llamada, el retardo algorítmico debido a ventanas o muestras futuras y la latencia de entrada a salida del dispositivo. Cumplir el plazo de un bloque no implica que la latencia total sea igual a $T_B$.

En el modelo neuronal se contabilizarán acumulación de tramas, salto, reconstrucción, adaptación entre tamaños de bloque y cualquier remuestreo. En ambos métodos se medirán percentiles y máximos observados de tiempo de procesamiento, además del número de incumplimientos. Las mediciones se realizarán sobre sesiones sostenidas y con el dispositivo de destino.

\subsection{Posicionamiento de la investigación}
Los dos antecedentes seleccionados motivan estudiar estrategias de identificación y de supresión aprendida dentro de un mismo montaje. La tesis no buscará establecer una clasificación universal entre familias de algoritmos. Su alcance será explicar el comportamiento de dos métodos concretos y de las variantes identificadas, bajo un protocolo compartido de estabilidad, calidad y ejecución.

\section{Metodología}
\subsection{Enfoque general y etapas de validación}
Se realizará una investigación experimental y comparativa con tres niveles de evidencia: simulación con señales de referencia conocidas, verificación numérica de implementaciones y medición física en lazo cerrado. Los resultados se reportarán por separado. La simulación permitirá controlar el camino y conocer el retorno verdadero; los ensayos físicos permitirán comprobar el comportamiento del sistema completo.

\subsection{Banco experimental y condiciones comunes}
Se utilizará un sistema mono compuesto por micrófono, interfaz de audio, computador y parlante de realimentación. Una fuente independiente reproducirá la voz de prueba cuando se necesite repetir exactamente el estímulo. Se registrarán modelo de equipos, ubicación, orientación, distancias, sala, ganancias analógicas y digitales, frecuencia y tamaño de bloque efectivos, controlador y versión del software.

Se trabajará a 16 y 48 kHz cuando el hardware y las implementaciones lo permitan. La comparación principal se efectuará a una tasa común validada para ambos métodos; las extensiones se reportarán aparte. Se ensayarán bloques de 64, 128 y 256 muestras. Se documentará cualquier agrupación en tramas o conversión de tasa, con su costo y retardo, sin asumir que ambos métodos requieren el mismo tamaño interno de trama.

La voz será la señal principal de evaluación. Ruido y música se utilizarán como pruebas complementarias de excitación, comportamiento tonal o generalización, sin extrapolar a música los resultados de métricas diseñadas para voz.

\propuestafigura{Montaje experimental}{Tomar una fotografía del montaje y elaborar una planta a escala aproximada con distancias y orientaciones. Distinguir la fuente de voz del parlante de realimentación. Incorporar una captura del banco JUCE como evidencia de instrumentación, con fecha y configuración visibles.}

\subsection{Implementación y verificación de ANIA}
Se utilizará la implementación C++ existente como punto de partida y se contrastará con una referencia independiente en Python de doble precisión. Se comprobarán correlaciones, solución del sistema, retardo, longitud efectiva del FIR y señal cancelada. Se incluirán caminos conocidos, señales de baja energía, estimaciones rechazadas y cambios del modelo durante la captura.

En los ensayos iniciales la calibración será manual y se aplicará una segunda sonda para validar el filtro antes de activarlo. Se medirán duración, nivel e interrupción asociada a la calibración. Frente a cambios del camino se distinguirán dos condiciones: filtro sin actualizar y recalibración controlada. Si se incorpora detección automática, se medirán falsas alarmas, omisiones y demora de detección antes de integrarla con la reinyección de ruido. Mientras esto no se complete, la variante se identificará como ANIA con calibración manual.

\subsection{Datos, entrenamiento y verificación de L3C-DeepMFC}
Se elaborará una especificación de reproducción que relacione cada componente publicado con su implementación: representación de entrada y objetivo, módulos recurrentes, estados, ventanas, reconstrucción, funciones de pérdida y ajuste fino en lazo cerrado. Se verificará la ejecución secuencial y se contabilizará cualquier anticipación temporal. Las decisiones que la publicación no permita resolver se documentarán como supuestos de reproducción.

Se seleccionarán corpus de voz y respuestas al impulso con permisos de uso adecuados y diversidad de hablantes y caminos acústicos. Se congelarán particiones de entrenamiento, validación y prueba sin compartir identidades de hablantes ni versiones del mismo registro entre particiones. También se reservarán caminos o salas para evaluar generalización. Las mezclas derivadas de una misma señal conservarán su partición de origen.

La selección de hiperparámetros y del modelo se realizará únicamente con entrenamiento y validación. Las grabaciones propias se reservarán para evaluación externa. Se registrarán semillas, versiones, configuración, costo de entrenamiento y criterio de selección. Se compararán las salidas de referencia y de inferencia integrada, incluyendo continuidad de estados entre bloques, inicio y final de sesión y diferencias numéricas tras la exportación.

Si la reproducción fiel no pudiera completarse, se informará el alcance logrado y se acordará una reformulación del objetivo antes de presentar una variante como sustituto. La existencia de código exploratorio no se considerará evidencia de entrenamiento ni de desempeño del método objetivo.

\subsection{Diseño experimental}
Se compararán tres condiciones principales: sin cancelación, ANIA y L3C-DeepMFC. Un esquema híbrido de cancelación clásica y supresión residual podrá explorarse una vez completada la comparación principal y se identificará como condición adicional.

El protocolo se organizará en caminos estacionarios y caminos con cambio controlado. Se propone comenzar con tres geometrías del conjunto micrófono--parlante y un cambio reproducible por desplazamiento u objeto cercano. Un piloto permitirá fijar distancias, niveles y tiempos. Para cada configuración se propone un mínimo de cinco repeticiones independientes, reiniciando el estado del sistema y equilibrando el orden de los métodos. Este número se revisará según la variabilidad del piloto; no constituye un cálculo de potencia estadística.

Los ensayos se emparejarán por estímulo, geometría y configuración del dispositivo. Se evaluará calidad a una ganancia común y también cerca del límite estable de cada método, identificando ambos análisis. Cuando las latencias difieran, se informará el resultado con latencia nativa y, si es viable, una condición con retardo total equiparado. No se interpretará como ventaja algorítmica una diferencia debida únicamente al retardo del lazo.

\subsection{Estabilidad y ganancia adicional}
La ganancia máxima estable (MSG) se definirá operacionalmente como la mayor ganancia que satisfaga el criterio de estabilidad durante una ventana de observación fija. Se propone un barrido de 1 dB con observación de 10 s por escalón, seguido de un ensayo sostenido en el nivel seleccionado. El piloto establecerá el umbral de persistencia tonal y crecimiento de energía para detectar inestabilidad; estos parámetros se fijarán antes de la campaña definitiva.

Se registrarán por separado acople, saturación y actuación del limitador. Si una protección o un límite del equipo impide alcanzar el umbral, la medición se reportará como acotada, no como una MSG exacta. La ganancia adicional estable será
\begin{equation}
\mathrm{ASG}=\mathrm{MSG}_{\text{método}}-\mathrm{MSG}_{\text{sin cancelación}},
\end{equation}
con ambas ganancias expresadas en dB y medidas bajo las mismas condiciones. Se registrará además el nivel de la señal útil para detectar estabilidad obtenida a costa de atenuación excesiva. Los ensayos usarán límites de salida y detención definidos durante el piloto.

\subsection{Calidad, error y robustez}
En simulación se calculará error de identificación cuando el camino verdadero sea conocido, y calidad de reconstrucción mediante SI-SDR\cite{leroux2019} y error respecto de la señal deseada alineada. Se conservará también una medida sensible al nivel, pues SI-SDR puede ocultar cambios de ganancia. El protocolo de alineamiento será común y no se utilizará para eliminar diferencias de latencia del reporte temporal.

Para voz se propone utilizar STOI\cite{taal2011} o ESTOI como indicadores de inteligibilidad y una métrica de calidad como PESQ o ViSQOL, según compatibilidad de tasa, disponibilidad e implementación validada. Se registrarán versión y modo de cada métrica. Las evaluaciones físicas con referencia requerirán una captura de la señal deseada en el micrófono con el lazo abierto y condiciones reproducibles; si esta referencia no es válida, no se presentarán esas puntuaciones como mediciones objetivas de calidad del montaje.

La reducción energética del retorno se calculará únicamente donde este pueda aislarse, por ejemplo en simulación o en ensayos de sonda. No se utilizará una razón entre entrada y salida del micrófono como ERLE cuando contenga voz deseada no separable. Ante cambios del camino se medirán pérdida de margen estable, degradación de calidad, demora de recuperación y tiempo dedicado a recalibración. Se registrarán asimismo fallos de recuperación.

\subsection{Ejecución en tiempo real y plataforma embebida}
Se medirán tiempos por bloque, percentiles 50, 95 y 99, máximo observado, incumplimientos del plazo, pérdidas de audio y uso de memoria y CPU. Se propone una prueba sostenida de al menos 10 minutos por configuración final, indicando carga del sistema y número de bloques. La latencia de entrada a salida se medirá mediante un procedimiento de retorno de señal; la propagación acústica se informará separadamente si se incluye en la medición.

La evaluación en Raspberry Pi 4 será posterior a la validación en PC e incluirá el mismo protocolo temporal, temperatura y posibles reducciones de frecuencia. Se documentarán las configuraciones que no puedan ejecutarse de forma sostenida. No se extrapolarán los tiempos del PC como resultados de la plataforma embebida.

\subsection{Análisis y reproducibilidad}
Se reportarán resultados por método y condición con medidas de dispersión e intervalos de confianza cuando el número de unidades independientes lo permita. Las diferencias se calcularán sobre ensayos emparejados. Para calidad se evitará tratar cada trama como observación independiente; se agrupará por grabación o hablante. Para estabilidad se considerarán las repeticiones y geometrías del montaje. Las conclusiones distinguirán magnitud del efecto, incertidumbre y relevancia práctica.

Cada sesión conservará audios, parámetros efectivos, eventos, modelo utilizado, marcas temporales y verificaciones de integridad. Se fijarán previamente las reglas de exclusión, como captura incompleta o saturación ajena al fenómeno estudiado, y se informará cuántos casos fueron excluidos. Los fallos propios de un método formarán parte de sus resultados. Se publicarán configuraciones y scripts reproducibles, respetando los permisos de los datos de voz.

\section{Alcance y Limitaciones}
El estudio se centrará en un sistema mono y en voz, con condiciones acústicas controladas y un número acotado de geometrías. No buscará validar una solución para todos los recintos ni establecer desempeño clínico en usuarios de audífonos. La transferencia desde los contextos de las publicaciones al banco de audio en vivo será una cuestión experimental.

Las conclusiones dependerán de la calidad de reproducción de los métodos, de la cobertura del corpus y de las restricciones del hardware. La evaluación embebida se limitará a Raspberry Pi 4. No se comprometerá el desarrollo de arquitecturas neuronales nuevas, arreglos multicanal ni una campaña perceptual con participantes como requisito de esta etapa. El control automático de ANIA y el esquema híbrido serán extensiones condicionadas a la validación de la comparación principal.

\section{Plan de trabajo}
La planificación se actualizará a partir del 28 de septiembre de 2026. Se propone una ejecución hasta marzo de 2027 para completar los hitos pendientes. Estas fechas constituyen una estimación de trabajo y deberán ajustarse con el profesor patrocinante al calendario académico y a la disponibilidad de equipos; no representan una fecha de entrega institucional confirmada.

\begin{table}[H]
\centering\small
\renewcommand{\arraystretch}{1.18}
\begin{tabular}{|p{0.06\textwidth}|p{0.49\textwidth}|p{0.16\textwidth}|p{0.16\textwidth}|}
\hline
\textbf{EDT}&\textbf{Actividad y producto verificable}&\textbf{Inicio}&\textbf{Término}\\\hline
1&Revisión del anteproyecto y especificación de métodos y protocolo&28-09-2026&09-10-2026\\\hline
2&Piloto físico, auditoría de capturas y congelación del protocolo&05-10-2026&23-10-2026\\\hline
3&Evaluación acústica de ANIA y documentación de variantes&19-10-2026&13-11-2026\\\hline
4&Selección del corpus, particiones y reproducción de L3C-DeepMFC&12-10-2026&20-11-2026\\\hline
5&Entrenamiento, ajuste fino y selección con validación&16-11-2026&18-12-2026\\\hline
6&Integración neuronal y equivalencia de inferencia continua&07-12-2026&08-01-2027\\\hline
7&Campaña comparativa física y evaluación externa&11-01-2027&05-02-2027\\\hline
8&Evaluación sostenida en Raspberry Pi 4&01-02-2027&19-02-2027\\\hline
9&Análisis, figuras, discusión y revisión de resultados&08-02-2027&05-03-2027\\\hline
10&Redacción y revisión de tesis&28-09-2026&19-03-2027\\\hline
\end{tabular}
\caption{Planificación propuesta de las actividades pendientes.}
\label{tab:plan}
\end{table}

La campaña comparativa comenzará cuando las implementaciones superen la verificación numérica y se hayan congelado parámetros y particiones. Si el entrenamiento o la integración requieren más tiempo, se actualizarán las etapas dependientes. La redacción y la documentación acompañarán todo el proceso.

\subsection{Carta Gantt resumida}
\begin{table}[H]
\centering\small
\renewcommand{\arraystretch}{1.2}
\begin{tabular}{|p{0.42\textwidth}|c|c|c|c|c|c|}
\hline
\textbf{Actividad}&\textbf{Oct.}&\textbf{Nov.}&\textbf{Dic.}&\textbf{Ene.}&\textbf{Feb.}&\textbf{Mar.}\\\hline
Protocolo y piloto&X&&&&&\\\hline
Evaluación ANIA&X&X&&&&\\\hline
Corpus y reproducción neuronal&X&X&&&&\\\hline
Entrenamiento y validación&&X&X&&&\\\hline
Integración continua&&&X&X&&\\\hline
Campaña comparativa&&&&X&X&\\\hline
Evaluación Raspberry Pi 4&&&&&X&\\\hline
Análisis y figuras&&&&&X&X\\\hline
Redacción y revisión&X&X&X&X&X&X\\\hline
\end{tabular}
\caption{Distribución mensual propuesta, octubre de 2026 a marzo de 2027.}
\end{table}

\section{Resultados esperados y avances preliminares}
\subsection{Resultados esperados}
Se espera obtener una caracterización comparativa de las regiones de operación de ambos métodos: ganancia utilizable, calidad conservada y costo de ejecución. Se elaborarán tablas de ganancia adicional estable, calidad de voz y respuesta a cambios del camino, junto con curvas de latencia y tiempo de procesamiento. Se identificarán las condiciones en que cada método resulte viable en PC y en Raspberry Pi 4.

Los productos incluirán las implementaciones documentadas, la especificación de reproducción neuronal, el protocolo, las configuraciones de los ensayos y los scripts de análisis. No se fija como resultado esperado que el aprendizaje profundo supere al método clásico; se buscará evidencia que permita aceptar o rechazar esa posibilidad en cada condición.

\subsection{Avances disponibles al 28 de septiembre de 2026}
Como antecedente de factibilidad, el proyecto dispone de un banco mono C++/JUCE con ganancia, sondas, estimación del camino, validación independiente y cancelación FIR con calibración manual. También existe un laboratorio Python para simulación y herramientas de entrenamiento y exportación de candidatos neuronales exploratorios.

Los registros internos del proyecto informan campañas sintéticas aprobadas de 50 casos y 25 pruebas matemáticas por frecuencia, a 16 y 48 kHz. La verificación de integración más reciente registra cinco pruebas C++ y 17 Python aprobadas, además de 12 sesiones sintéticas completas auditadas y dos sesiones incompletas correctamente rechazadas\cite{estado2026}. Estos antecedentes corresponden a verificaciones del software y a equivalencia numérica dentro de los casos ensayados.

Permanecen pendientes la campaña acústica física, la reproducción fiel y entrenamiento definitivo de L3C-DeepMFC, la integración neuronal y las mediciones sostenidas en Raspberry Pi 4. Tampoco se ha completado el control automático descrito para ANIA. Por tanto, los avances no constituyen una comparación final ni acreditan todavía ganancia acústica adicional, calidad perceptual o latencia física.

\propuestafigura{Evidencia preliminar de identificación}{Generar una gráfica con camino sintético conocido, estimación Python y estimación C++, acompañada de error y configuración. Identificarla explícitamente como validación numérica. Más adelante, agregar una respuesta al impulso medida, separada del ejemplo sintético.}

\propuestafigura{Resultados de la campaña futura}{Preparar gráficos de ASG por geometría, calidad frente a ganancia y percentil 99 de procesamiento frente al plazo del bloque. Para cambios de camino, usar una línea temporal con cambio, degradación y recuperación. No llenar estos gráficos con valores ilustrativos que puedan confundirse con resultados.}

\section{Criterios para la discusión de resultados}
La discusión relacionará las diferencias observadas con la forma de operación de cada método. Para ANIA se examinarán la precisión del camino identificado y el costo de recalibrar; para L3C-DeepMFC, la generalización a caminos no vistos y la continuidad de inferencia. Se distinguirá entre limitaciones del método, diferencias de reproducción y restricciones de la plataforma.

Se evaluará si una mejora de estabilidad mantiene la señal útil y si una mejora de calidad persiste al considerar el retardo total. Las condiciones con fallos, resultados negativos o inviabilidad temporal se incorporarán al análisis. Los ensayos complementarios no reemplazarán la comparación principal y se identificarán como exploratorios.

\section{Síntesis de la propuesta}
Se propone comparar ANIA y L3C-DeepMFC en un banco de realimentación acústica en tiempo real, mediante verificación numérica, simulación y evaluación física. La investigación buscará determinar los compromisos entre estabilidad, conservación de voz, robustez y recursos computacionales, con trazabilidad de datos y configuraciones.

La infraestructura existente permitirá iniciar el piloto físico y desarrollar la reproducción neuronal, pero las conclusiones sobre desempeño se establecerán después de completar la campaña. El resultado de la tesis será una evaluación fundamentada de las condiciones de uso de ambos métodos y de su factibilidad en una plataforma embebida.

\renewcommand{\refname}{Bibliografía}
\addcontentsline{toc}{section}{Bibliografía}
\begin{thebibliography}{9}
\bibitem{patel2020}
K. Patel e I. M. S. Panahi, ``Efficient Real-Time Acoustic Feedback Cancellation using Adaptive Noise Injection Algorithm'', \textit{42nd Annual International Conference of the IEEE Engineering in Medicine \& Biology Society (EMBC)}, pp. 972--975, 2020.
\url{https://doi.org/10.1109/EMBC44109.2020.9175686}.
Texto disponible en \url{https://pmc.ncbi.nlm.nih.gov/articles/PMC7545262/}.

\bibitem{hao2025}
F. Hao, B. C. J. Moore, H. Zhang, X. Li y C. Zheng, ``L3C-DeepMFC: Low-Latency Low-Complexity Deep Marginal Feedback Cancellation with Closed-Loop Fine Tuning for Hearing Aids'', \textit{Interspeech 2025}, pp. 838--842, 2025.
\url{https://doi.org/10.21437/Interspeech.2025-1111}.
Texto disponible en \url{https://www.isca-archive.org/interspeech_2025/hao25_interspeech.pdf}.

\bibitem{leroux2019}
J. Le Roux, S. Wisdom, H. Erdogan y J. R. Hershey, ``SDR -- Half-Baked or Well Done?'', \textit{IEEE International Conference on Acoustics, Speech and Signal Processing (ICASSP)}, 2019.
\url{https://doi.org/10.1109/ICASSP.2019.8683855}.

\bibitem{taal2011}
C. H. Taal, R. C. Hendriks, R. Heusdens y J. Jensen, ``An Algorithm for Intelligibility Prediction of Time-Frequency Weighted Noisy Speech'', \textit{IEEE Transactions on Audio, Speech, and Language Processing}, vol. 19, n.\ 7, pp. 2125--2136, 2011.
Texto de los autores: \url{https://ceestaal.nl/Taal(2011).pdf}.

\bibitem{estado2026}
M. I. Arancibia Villanueva, ``Estado de la tesis y registros de verificación del banco experimental'', documentación interna del proyecto \textit{ANIA-vs-L3C-DeepMFC}, actualización del 28 de septiembre de 2026. Archivos \texttt{IMPLEMENTATION\_STATUS.md}, informes de la referencia independiente ANIA y registro de integración de captura del 28 de septiembre de 2026. Evidencia preliminar no publicada.
\end{thebibliography}


\end{document}
